\section*{6.3 Deterministic communication lifting support (author source pp.41--45)}
\begin{originalresult}{Theorem 6.9}
If a two-party gadget $g:X\times Y\to\bits$ satisfies $\AND,\XOR\le g$, then for every function $f:\bits^n\to\bits$, the function $f\comp g$ has communication complexity $\Omega(\operatorname{bs}(f)^k)$ for some fixed constant $k>0$. The same is true if $g$ satisfies $\OR,\XOR\le g$.
\end{originalresult}
We start with the following observation: according to Hatami et al.(2018), for any $f:\bits^n\to\bits$ we have $\DT_\oplus(f)\le O(D(f\comp\XOR)^6)$. Zhang \& Shi(2010) showed that $\DT_\oplus(f)=O(C_\oplus(f)^2)$. Combining these results, we can derive the following corollary:
\begin{originalresult}{Corollary 6.11}
For any $f:\bits^n\to\bits$, we have
\[C_\oplus(f)\le O(D(f\comp\XOR)^{12}).\]
\end{originalresult}
Using this corollary, our main goal will be to find a substitution $\tau$ and two points $x$ and $y$ such that $f(x)\ne f(y)$ and $x|_\tau$ and $y|_\tau$ have parity certificates of large minimal weight. Then, using Lemma6.12 and Lemma6.13, we will get the required lifting.
In the following lemma, a minimal sensitive block of a function $f:\bits^n\to\bits$ at point $x$ is a subset $I\subseteq[n]$ such that $f(x^I)\ne f(x)$ but $f(x^J)=f(x)$ for all $J\subsetneq I$. This lemma shows that the deterministic communication complexity of $f\comp\XOR$ is lower bounded by the size of any minimal sensitive block of $f$.
\begin{originalresult}{Lemma 6.12}
For any $f:\bits^n\to\bits$ and $\alpha\in\bits^n$ such that there is a minimal sensitive block of $f$ at $\alpha$ of size $S$, the following holds:
\[D(f\comp\XOR)\ge S.\]
\end{originalresult}
\begin{proof}
Consider some minimal sensitive block of $f$ at $\alpha$ of size $S$. Let $i_1,i_2,\ldots,i_S$ be the indices of the variables in this block of $\alpha$. Let us restrict the function $f$ by substituting values $x_j=\alpha_j$ for all $j\notin\{i_1,i_2,\ldots,i_S\}$. This gives us a function $f'(y_1,y_2,\ldots,y_S)$ such that
\begin{itemize}
\item $f'(y_1,\ldots,y_S)=f(\alpha)$, if $y_j=\alpha_{i_j}$ for at least one $j$.
\item $f'(y_1,\ldots,y_S)\ne f(\alpha)$, if $y_j\ne\alpha_{i_j}$ for all $j$.
\end{itemize}
Thus, if we substitute XOR inside each $y_j$ this will give us an instance of the function
\[(z_1\oplus z_1'\oplus\alpha_{i_1})\vee(z_2\oplus z_2'\oplus\alpha_{i_2})\vee\cdots\vee(z_S\oplus z_S'\oplus\alpha_{i_S}),\]
where each $z_i$ and $z_i'$ are distinct variables. This function is essentially a negation of the equality function with a shift. So, its deterministic complexity equals $S$ (see Kushilevitz \& Nisan(1996) for the lower bound for $\operatorname{EQ}_n$). Thus,
\[D(f\comp\XOR)\ge D(f'\comp\XOR)\ge S.\]
\end{proof}
The following lemma shows that the deterministic communication complexity of $f\comp\AND$ is lower bounded by the block sensitivity of $f$ at the all-0 input.
\begin{originalresult}{Lemma 6.13}
For any $f:\bits^n\to\bits$ such that $\operatorname{bs}(f,0)=S$, the following holds:
\[D(f\comp\AND)\ge\Omega(S).\]
\end{originalresult}
\begin{proof}
The proof of the lemma essentially follows the proof from Zhang(2009). Let $I_1,\ldots,I_S$ be the minimal sensitive blocks of $f$ at 0. We describe a reduction from $\operatorname{PromiseDisj}_S$ to $f\comp\AND$, so any protocol for $f\comp\AND$ can be used to solve $\operatorname{PromiseDisj}_S$, and thus $D(f\comp\AND)\ge\Omega(S)$. For any input $(a,b)\in\bits^S$ to $\operatorname{PromiseDisj}_S$ we define the corresponding input $(x,y)\in\bits^n$ to $f\comp\AND$.
\begin{enumerate}
\item Let $x_i=y_i=0$ for all $i\notin\bigcup_{j=1}^S I_j$.
\item Let $x_i=a_j$ and $y_i=b_j$ for all $i\in I_j$.
\end{enumerate}
We claim that this is a correct reduction. Indeed, if there is a unique $j$ such that $a_j\wedge b_j=1$, then
\[f(x_1\wedge y_1,x_2\wedge y_2,\ldots,x_n\wedge y_n)=f(0^{I_j})\ne f(0,\ldots,0).\]
On the other hand, if $a_j\wedge b_j=0$ for all $j\in[S]$, then
\[f(x_1\wedge y_1,x_2\wedge y_2,\ldots,x_n\wedge y_n)=f(0,\ldots,0).\]
\end{proof}
The next lemma is an application of the `minimum weight lemma' for the identity gadget.
\begin{originalresult}{Lemma 6.14}
Suppose that $f$ has a parity certificate at some point $x$ of size at most $k$. Let $K$ be a parameter. Suppose that the certificate complexity of $f$ at $x$ is greater than $k\cdot K$. Then we can find a partial assignment $\sigma$ to at most $k\cdot K$ variables of $f$, consistent with $x$, such that:
\begin{itemize}
\item $f|_\sigma$ has a parity certificate at $x|_\sigma$ of size at most $k$.
\item The minimum weight of this certificate is at least $K$.
\end{itemize}
\end{originalresult}
\begin{proof}
Apply Lemma5.3 for the function $f$ and the identity function as a gadget $g$.
\end{proof}
For the last ingredient needed for the proof of the classification theorem, the following lemma shows that parity certificates of high minimum weight always intersect with (regular) certificates of sufficiently small size.
\begin{originalresult}{Lemma 6.15}
Suppose that for some function $f$ and some input $x$ the following holds: $f$ has a parity certificate $C_\oplus(x)$ of size $k$ at $x$ whose minimum weight is at least $K$. Then, for any point $y$ with certificate complexity less than $K$, there is a non-empty intersection of the corresponding certificates $C(y)$ and $C_\oplus(x)$.
\end{originalresult}
\begin{proof}
Apply Lemma5.4 with $C_\oplus=C_\oplus(x)$ and $C=C(y)$.
\end{proof}
Now we have everything we need to prove Theorem6.9.
\begin{proof}[Theorem 6.9]
Assume that $\AND,\XOR\le g$. The proof in the case of $\OR,\XOR\le g$ is similar. Consider some function $f$. There is a value $S$ such that $\operatorname{bs}(f)=9S^5$. If $C_\oplus(f)\ge S$, then by Corollary6.11,
\[D(f\comp g)\ge D(f\comp\XOR)\ge C_\oplus(f)^{1/12}\ge S^{1/12}\ge\Omega(\operatorname{bs}(f)^{1/(12\cdot5)}),\]
and so the theorem holds for $k=1/60$. So, we can assume that $C_\oplus(f)<S$. Also, we can assume that for any input, the size of every minimal sensitive block is at most $S$, otherwise by Lemma6.12 we get
\[D(f\comp g)\ge D(f\comp\XOR)=\Omega(S)=\Omega(\operatorname{bs}(f)^{1/5}).\]
Consider any input $x$ with maximal block sensitivity and the corresponding set of minimal sensitive blocks. The certificate complexity of $x$ is at least $9S^5$, since the certificate complexity is at least the block sensitivity. By Lemma6.14 (for $k=S$ and $K=2S^3$), we can find a substitution $\sigma$ such that $f|_\sigma$ has a parity certificate at $x|_\sigma$ of size at most $S$ with minimal weight at least $2S^3$. In addition, the block sensitivity of $f|_\sigma$ at $x|_\sigma$ is at least $9S^5-2S^4\ge7S^5$ (since we set at most $2S^4$ variables with $\sigma$), and the blocks still have size at most $S$. We fix the substitution $\sigma$ and consider the restricted function. In what follows, we denote $\hat f:=f|_\sigma$ and $\hat x:=x|_\sigma$.
Consider any input $y$ obtained from $\hat x$ by flipping one of the minimal sensitive blocks. Suppose that the certificate complexity of this input is at most $2S^3-S$. Then by Lemma6.15 (for $k=C_\oplus(\hat f,\hat x)\le S$ and $K=2S^3$), the certificate for $y$ and the parity certificate for $\hat x$ intersect. Since $\hat f(\hat x)\ne\hat f(y)$, this leads to a contradiction. So, the certificate complexity of $y$ should be greater than $2S^3-S$.
Again, by using Lemma6.14 (for $k=S$ and $K=S^2-1$), there is a substitution $\tau$ of at most $S^3$ variables, such that $y|_\tau$ has a parity certificate of size at most $S$ with minimal weight at least $S^2-1$. In addition, the certificate complexity of $y|_\tau$ is at least $2S^3-S^3=S^3$.
Since the substitution $\tau$ consists of at most $S^3$ variables, the minimal weight of the parity certificate $C_\oplus(\hat f,\hat x)|_\tau$ (we apply $\tau$ to every linear form in the certificate) is at least $2S^3-S^3\ge S^3$. Note that substitution $\tau$ can be inconsistent with $\hat x$, but we still can make this substitution inside $C_\oplus(\hat f,\hat x)$ and get a solvable certificate for $\hat f|_\tau$. Thus, if we find any point $z$, consistent with the substitution $\tau$, such that the certificate complexity of $z|_\tau$ is at most $S^2-S-1$, then this certificate $C(\hat f|_\tau,z|_\tau)$ should intersect with both $C_\oplus(\hat f,\hat x)|_\tau$ and $C_\oplus(\hat f,y)|_\tau$ by Lemma6.15 (for $k=C_\oplus(\hat x)\le S$ and $K=S^2-1$). Since all the these certificates intersect, we have $\hat f(\hat x)=\hat f(z)=\hat f(y)$, but this contradicts $\hat f(\hat x)\ne\hat f(y)$.
Now, consider an input $z$ consistent with $\tau$ such that all the $*$-variables of $\tau$ are set to zeroes. We have shown above that the certificate complexity of $\hat f|_\tau$ at $z|_\tau$ is at least $S^2-S$. Recall that $z|_\tau$ has a block sensitivity witness with all the blocks of size at most $S$. Since all the blocks together constitute a certificate, we conclude that the block sensitivity of $\hat f|_\tau$ at $z|_\tau$ is at least $S-1$. Since $z|_\tau$ is the all zeroes, by Lemma6.13 we have
\[D(f\comp g)\ge D(f\comp\AND)=\Omega(S)=\Omega(\operatorname{bs}(f)^{1/5}).\]
\end{proof}
